% Appendix section: executor design search on GameBoyWorlds (negative result), contrasted with WebVoyager.
% Requires \usepackage{booktabs} (and \usepackage{multirow} is not needed).

\section{Executor design search on GameBoyWorlds}
\label{app:gameboy-executor-search}

A natural objection to our GameBoyWorlds results is that the executor we use is too weak, and that a better-designed
acting agent, with the same model, would close much of the gap. To test this, we ran a systematic, pre-registered
search over executor designs, with the model held fixed. We ran exactly the same search on WebVoyager, where it found
changes that matter: a $+10.0$ point gain on the full benchmark that holds on held-out tasks. On GameBoyWorlds, the
same procedure screened 18 executor features over more than 3{,}000 dev episodes. None of them improved success: the
best two moved it by $+1.2$ points, and 16 of 18 scored \emph{below} the baseline. This section documents the search
so that the negative result can be judged on its own terms.

\subsection{Setup}

\paragraph{Fixed parts.} One open-weight vision-language model (Gemma-4-31B-it) is used for every model call; it is
not fine-tuned, and sampling is unseeded, so repeats are independent runs. The supervisor of our agent framework is
switched off in every arm, so all differences come from the executor. The baseline executor is the benchmark's native
agent: at each step it sees the screenshot and its recent action history and presses one button. The step budget
(175) and the success checks are the benchmark's own. The agent may not read game memory or the ROM; the one
exception, used by a single feature (F04), is a reader for the text that the console is displaying on screen. There is
no save-state lookahead.

\paragraph{Splits.} Before any feature work, the 99 tasks (Pok\'emon Red and Crystal) were split once into 50 dev
tasks (25 per game) and 49 held-out tasks. All selection used dev only.

\paragraph{Noise floor.} The baseline was run 5 times on dev: 34, 32, 34, 38 and 36\% (mean 34.8\%). The floor is
the largest gap between two baseline repeats or $1.96\cdot\mathrm{SD}(\text{paired diff})/\sqrt{n}$, whichever is
larger: 10.5 points for a single repeat. Single runs are therefore far too noisy to judge a feature on this
benchmark, so every feature ran at least 3 times, giving a screening threshold of $10.5/\sqrt{3} = 6.0$ points at 3
repeats and $5.25$ at 4.

\paragraph{Statistics.} Arms are compared task by task on the same tasks. Success is the mean over tasks of the
per-task mean over repeats. 95\% confidence intervals bootstrap over tasks (5{,}000 resamples); W/L counts the tasks
an arm wins or loses against the baseline.

\subsection{Procedure}

The procedure was written down before each stage was run, and is identical across benchmarks except for the
thresholds, which follow each benchmark's noise floor.
\begin{enumerate}
  \item \textbf{Catalogue.} A survey of agent papers and open-source harnesses produced 54 executor-side features,
  each with its prior evidence and a per-benchmark implementation note, ranked per benchmark by evidence over cost.
  \item \textbf{Failure taxonomy.} 50 failed baseline dev episodes were labelled with one main failure each and
  spot-checked against the trajectories: loop or no-op 27, screen misread 15, wrong plan 4, out of steps 4. The
  shortlist was checked to cover the two dominant buckets.
  \item \textbf{Smoke test.} Each prototype first ran live on 4 tasks, and the episode records were read to confirm
  that it did what it claimed.
  \item \textbf{Screening.} Each feature alone against the baseline on the 50 dev tasks. \emph{Promote} if the gain
  is at least the threshold at under twice the baseline's model calls, or if it is neutral but clearly fixes its
  target failure bucket; \emph{drop} if the gain is at most minus the threshold, or the calls double for no gain;
  otherwise \emph{retest} once with one more repeat if the feature is in the top half of the prior ranking, and drop
  it if not.
  \item \textbf{Combination, leave-one-out and confirmation.} Promoted features would be combined greedily, pruned by
  leave-one-out, and the final design confirmed on held-out tasks. On GameBoyWorlds no feature was promoted, so these
  stages did not run.
\end{enumerate}

\subsection{Features screened}

Table~\ref{tab:gb-screening} lists every feature screened. They span every family in the catalogue that has a
GameBoyWorlds counterpart:
\begin{itemize}
  \item \textbf{Perception:} on-screen text as a text block (F04), $3\times$ screenshot upscaling (F07), and game-mode
  detection (overworld, dialogue, menu, battle) from the screen (F54).
  \item \textbf{Memory and history:} a longer action history (F26), a summarised history (F25), and an objectives
  and places log written by the model (F28).
  \item \textbf{Planning and knowledge:} a persistent plan (F21), a tips document (F29) and a game guide (F30), both
  written from the model's prior knowledge.
  \item \textbf{Loop handling:} a no-op ban that forbids repeating a press that changed nothing (F31), and a loop
  warning when the same screen recurs (F32).
  \item \textbf{Reasoning and verification:} a structured reasoning template (F19), a pre-execution critic (F38),
  and a self-consistency vote over three samples (F37).
  \item \textbf{Action space and budget:} several presses per decision (F16), macro actions such as ``walk $n$
  tiles'' and ``advance dialogue'' (F17), and the remaining step budget shown each step (F10).
  \item \textbf{An expectation line} (F62), taken from a recent Android agent: the model states the screen change it
  expects from its action, and the next prompt shows that expectation and asks the model to compare it with the new
  screen. We added this one late, as the only mechanism of two scaffolds that report solving AndroidWorld that was
  not already covered here. It targets both main failure buckets (the model must notice when a press did not do what
  it believed) at no extra model call.
\end{itemize}
Between them, the features aim at both dominant failure buckets: loops (F31, F32, F38, F62, F10, F21) and screen
misreads (F04, F07, F54, F62).

\begin{table}[t]
\centering
\small
\caption{Executor features screened on the 50 GameBoyWorlds dev tasks, each alone against the baseline. The
threshold is 6.0 points at 3 repeats and 5.25 at 4. Calls are model calls per task. No feature was promoted.}
\label{tab:gb-screening}
\begin{tabular}{llrrrrl}
\toprule
ID & Feature & Rep. & Success & $\Delta$ vs base [95\% CI] & W/L & Calls \\
\midrule
-- & Baseline (native executor) & 5 & 34.8 & -- & -- & 228 \\
\midrule
F07 & Screenshot upscaling ($3\times$) & 4 & 36.0 & $+1.2$ [$-5.3$, $+8.0$] & 8/7 & 231 \\
F62 & Expectation line & 4 & 36.0 & $+1.2$ [$-4.8$, $+6.9$] & 11/8 & 233 \\
F32 & Loop warning & 3 & 34.7 & $-0.1$ [$-5.6$, $+6.0$] & 6/7 & 231 \\
F31 & No-op ban & 3 & 34.0 & $-0.8$ [$-5.9$, $+3.7$] & 7/6 & 242 \\
F38 & Pre-execution critic & 3 & 33.3 & $-1.5$ [$-8.9$, $+6.4$] & 6/8 & 394 \\
F28 & Objectives and places log & 3 & 32.7 & $-2.1$ [$-11.3$, $+6.5$] & 7/11 & 239 \\
F10 & Step budget shown & 3 & 32.0 & $-2.8$ [$-7.1$, $+1.2$] & 3/6 & 234 \\
F16 & Several presses per decision & 3 & 32.0 & $-2.8$ [$-10.3$, $+4.4$] & 8/11 & 180 \\
F29 & Tips document from priors & 3 & 32.0 & $-2.8$ [$-11.9$, $+6.4$] & 6/7 & 225 \\
F04 & On-screen text as text & 4 & 31.0 & $-3.8$ [$-8.5$, $+0.2$] & 3/8 & 238 \\
F54 & Game-mode detection & 4 & 31.0 & $-3.8$ [$-8.4$, $+0.4$] & 3/9 & 235 \\
F37 & Self-consistency vote ($k{=}3$) & 3 & 30.7 & $-4.1$ [$-9.6$, $+0.9$] & 4/8 & 492 \\
F19 & Structured reasoning template & 3 & 30.7 & $-4.1$ [$-10.3$, $+1.9$] & 3/8 & 242 \\
F17 & Macro actions & 4 & 29.5 & $-5.3$ [$-13.2$, $+1.8$] & 5/8 & 198 \\
F26 & Longer history (15 steps) & 3 & 29.3 & $-5.5$ [$-10.1$, $-1.5$] & 2/8 & 241 \\
F25 & Summarised history & 3 & 29.3 & $-5.5$ [$-13.3$, $+1.5$] & 4/8 & 120 \\
F21 & Persistent plan & 3 & 28.0 & $-6.8$ [$-14.3$, $-0.0$] & 4/14 & 250 \\
F30 & Game guide from priors & 2 & 22.0 & $-12.8$ [$-21.2$, $-5.6$] & 1/11 & 262 \\
\bottomrule
\end{tabular}
\end{table}

\subsection{Results}

\paragraph{No feature helps.} The two best features gain $+1.2$ points, less than a quarter of the threshold; every
other feature scores at or below the baseline. Four borderline features with a high prior rank (F07, F04, F54, F17) were
retested with a fourth repeat, and none came near the threshold. Three features are
reliably \emph{harmful}, with confidence intervals at or below zero: a longer history (F26), a persistent plan (F21)
and a game guide written from the model's priors (F30). The last was stopped after 2 of 3 repeats by an early-drop
rule we added midway through screening to free compute; the rule was applied only to this arm. Since nothing was
promoted, the final GameBoyWorlds executor is the baseline. On the full benchmark it scores 31.3\% (99 tasks):
38.0\% on the dev half and 24.5\% on the held-out half.

\paragraph{The search could have detected a meaningful gain.} With 3--5 repeats per arm, each confidence interval
spans roughly $\pm 6$ points. No feature's upper bound exceeds $+8.0$ points, and only the two $+1.2$ features have
upper bounds above $+6.5$. So the search rules out any single feature worth more than about 8 points, which is less
than the $+10$ points the same search found on WebVoyager. The negative result is not a matter of too few runs: the
screening used 63 runs of the 50 dev tasks (more than 3{,}000 episodes) and more repeats per arm than on WebVoyager,
because the benchmark is noisier.

\paragraph{The failures the features target did not move.} Loops and screen misreads account for 42 of 50 labelled
failures. Neither loop handling (F31, F32, F38, F62) nor richer screen inputs (F04, F07, F54) moved either bucket.
Loop handling was hard even to target. An offline analysis of 150 baseline episodes found that revisiting a screen is
normal for this agent: the loop warning fired on 35\% of all action calls. Neither exact-frame revisits nor image
embedding similarity separated loop failures from other failures (AUC 0.50--0.58).

\paragraph{Reading partial runs would have produced a false positive.} While F62 was running, the table at 35 of
the 50 tasks read $+7.4$ points, above the threshold. The tasks still running were the long ones, which mostly fail;
with all 50 tasks, the gain was $+1.2$. All decisions above were taken only on complete runs.

\subsection{Contrast: the same search on WebVoyager}
\label{app:gameboy-vs-web}

The same procedure, with the same model, the same catalogue and the same rules, did find improvements on WebVoyager
(250 dev tasks; threshold 4.8 points at 1 repeat). Of 18 arms screened (15 features, two of them in several
variants), 4 were promoted at screening, all with confidence intervals above zero:
\begin{itemize}
  \item the remaining step budget shown each step (F10), $+7.2$ [$+3.3$, $+11.1$];
  \item a forced answer near the budget end (F11), $+6.4$;
  \item removing the agent's \texttt{Google} action, which often led to a CAPTCHA page the agent could not leave
  (F14), $+5.6$;
  \item five screenshots of history instead of three (F08), $+5.2$.
\end{itemize}
Greedy combination and leave-one-out kept F10, F14 and F08; each costs at least 2.7 points when removed. The final
design raises success from 56.4\% to 66.4\% on the full benchmark ($+10.0$ [$+5.8$, $+14.2$], 429 tasks) and from
57.0\% to 66.5\% on the 179 held-out tasks ($+9.5$ [$+2.8$, $+16.2$]), with no extra model calls. The search
procedure can therefore find real, replicable executor gains when they exist.

Table~\ref{tab:gb-vs-web} compares the features screened on both benchmarks. Features that helped on WebVoyager were
flat or harmful on GameBoyWorlds. Showing the step budget is $+7.2$ on the web and $-2.8$ on GameBoyWorlds, and more
visual history is $+5.2$ on the web while a longer history is $-5.5$ on GameBoyWorlds. The reasoning add-ons that
were flat on the web (plans, templates, votes, tips) were clearly negative on GameBoyWorlds.

\begin{table}[t]
\centering
\small
\caption{The same features screened on WebVoyager (250 dev tasks) and on GameBoyWorlds (50 dev tasks), as gains in
points over each benchmark's baseline. Bold marks a feature promoted at screening.}
\label{tab:gb-vs-web}
\begin{tabular}{lrr}
\toprule
Feature & WebVoyager & GameBoyWorlds \\
\midrule
Step budget shown (F10) & $\mathbf{+7.2}$ & $-2.8$ \\
More history (F08 five screenshots; F26 15 steps) & $\mathbf{+5.2}$ & $-5.5$ \\
On-screen or page text as text (F04) & $+2.8$ & $-3.8$ \\
Tips document from priors (F29) & $+2.4$ & $-2.8$ \\
Loop warning (F32) & $+1.6$ & $-0.1$ \\
Persistent plan (F21) & $+0.8$ & $-6.8$ \\
Self-consistency vote, $k{=}3$ (F37) & $+0.8$ & $-4.1$ \\
Structured reasoning template (F19) & $-1.2$ & $-4.1$ \\
\midrule
Arms screened & 18 & 18 \\
Promoted at screening & 4 & 0 \\
Final design vs baseline, held-out & $+9.5$ [$+2.8$, $+16.2$] & -- (baseline kept) \\
\bottomrule
\end{tabular}
\end{table}

\paragraph{Why the web gains do not transfer.} The WebVoyager gains came from the \emph{harness}, not from extra
reasoning. Showing the budget made the agent answer before running out of iterations. Removing a trap action
fixed the single largest failure cause in the web failure taxonomy (22 of 50 failures). More screenshots helped the
agent keep track of multi-page navigation. GameBoyWorlds has no comparable harness defect. Its budget is long (the
baseline uses 116 of 175 steps on average), its action space is a handful of console buttons, and its failures are
loops and misreads of a low-resolution game screen. Prompting changes do not fix those for this model.

\subsection{Limitations}

Three shortlisted GameBoyWorlds features with high build cost were not built: a labelled grid overlay on the
screenshot, a tile map of the screen built from the model's own per-tile classification, and a navigation tool that
walks to a chosen cell over that map. Such tools appear in hobbyist Pok\'emon-playing harnesses with frontier models;
we did not build them because each needs substantial engineering, and the tile map adds model calls. Thinking mode was not tested on any benchmark. Within these limits,
the evidence is that, for this model, no prompting-level or harness-level change to the executor that we could find
improves GameBoyWorlds performance. A meaningful improvement on this benchmark appears to need more than executor
design.
